% =====================================================================
%  R. Nielsen, OM SELVFORSTAAELSE I DET ONDE, MED HENBLIK PAA SHAKSPEARES RICHARD DEN TREDIE
%  In: For Ide og Virkelighed. Et Tidsskrift, udgivet af R. Nielsen, Bj. Bjørnson og Rud. Schmidt.  Aargang 1869.
%  Kjøbenhavn: Chr. Steen & Søns Forlag; Louis Kleins Bogtrykkeri.  pp. 521-554.  Antiqua.
%
%  TRANSCRIPTION -- GENERATED; DO NOT EDIT BY HAND.  The five Nielsen essays of the 1869 volume are built by ONE pipeline
%  (ebook-method/for-idee-1869/) into a combined working file, which split5.py cuts into the five transcriptions.
%  Rebuild from the repo root:   sh ebook-method/for-idee-1869/rebuild.sh
%
%  STATUS: COMPLETE (2026-10-13): the volume's title page and the whole essay, pp. 521-554, with its notes.
%
%  EDITION.  The only printing.  Text of record: the Internet Archive's scan of the University of Toronto copy
%  (forideogvirkelig00kjobuoft), ~/bibliotek/Nielsen, Rasmus/1869-for-idee-og-virkelighed.pdf (SCANS.tsv), 886 PDF pages,
%  400 ppi colour; text layer = the IA's ABBYY FineReader OCR.  Printed page N = PDF N + 10 to p. 688; a plate is bound in
%  before p. 689, so pp. 689-706 = PDF N + 12.  Title page PDF 7.  No e-book, no other scan, no earlier transcription found
%  (checked 2026-10-13: the repo, the web, the Internet Archive; HathiTrust and Google Books could not be queried that day).
%
%  ACCURACY (2026-10-13; one collation for all five essays, 189 pp., ~48,000 words).  The words are those of the ABBYY layer
%  (draft.py), collated word by word with an independent tesseract reading of the images (collate/check.py, witness 'ocr'):
%  1132 disagreements at the first collation.  check.py's own rules settled 96; rules.py settled ~330 more, each rule
%  switched on only after a BLIND PILOT of 107 crops (reader A, 21/22 controls) had measured it on this book: tesseract's
%  glyph faults 8/8, aa read an 8/8, accents it adds 8/8, the layer's glyph faults 7/7 and 8/8, o/ø by the ø rule 8/8, the
%  printer's space before a mark 8/8 (house rule), and the narrow word-division rules; the quote, dash, note-number and
%  word-seam rules failed the pilot and stayed off.  The rest were read at the image BLIND (sheets.py: 579 word crops on
%  58 sheets, two planted controls a sheet, 116 of 116 controls right; six readers) and about 110 spots in the main
%  session from line crops (zoom.py).  STRUCTURE read at the image (reader K): the essays' headings, the volume's title
%  page, every footnote whole (the notes were collated where they stand and then moved to their marks), the signatures.
%  LETTERSPACING (the print's emphasis, \emph): candidate lines by box width (spacedwidth.py; the gap measure is noise on
%  this scan): 163 lines read blind (22/22 roman controls), doubtful ones read in the main session.  ITALICS by stroke
%  slant (check.py --italics): 24 crops read blind; the only italics are Latin quotations.
%  SWEEPS after OPEN 0, each spot read at the image: the W-shortens / T-longer list of the blind decisions (found two
%  W decisions that had dropped a note's mark and a closing «, and one that had joined „der f.“); --moved and --splits;
%  the split sweep; the reader-line sweep; paragraph starts against tesseract's indents; quotation marks (the print
%  closes with « and often with », both as printed); the unknown-word sweep (31 spots: 16 corrections).
%  TRANSLATION as a last reading: the twelve translators listed 79 odd spots (ebook-method/for-idee-1869/tr/odd-*.tsv),
%  ~60 read at the image (zoom zo, zw, zx): mended -- f's whose head did not ink, read t by both OCRs (pp. 21, 36, 309, 319,
%  441; a recurring defective sort, as in Folkelige Foredrag), two blind decisions that had taken a reader's t for the
%  print's f (pp. 322, 539), a speck read as a stop (p. 589), a lost dash (p. 525), the verse on p. 414 set as one quotation,
%  false paragraph breaks at page turns left by the notes' removal (pp. 430, 553, 620, 695, 703); the rest are the print's.
%  SAMPLED PAGE CHECK (pagecheck.py, seed 202610136): pp. 12, 418, 521, 694 and the title page read whole and blind, 1069
%  words: 0 word and 0 punctuation errors; one letterspaced passage in a note missed (p. 694), mended.
%  What this cannot see: a misprint both OCRs normalise into a real word.
%  The ø rule (Hans 2026-10-08): o-forms that are no word (Oieblik, Ifolge, imode ...) are set with ø.
%  PRINTED AS IS p. 541: „Over-|enstemmelse“ with one s
%  Not transcribed: running heads, folios, sheet signatures („For Ide og Virkelighed.  20“), the library's stamp; the large
%  initial of each essay is set as a plain letter.  Quotation marks as printed (» « in the text, „ “ in the notes).  The notes
%  are numbered ¹) per page, as printed (p. 553: *)).
%
%  TOOLS (ebook-method/for-idee-1869/, copied from folkelige-foredrag): draft.py, ocr.py, rules.py, sheets.py + reads.py +
%  READER-BRIEF.md, heads.py, spacedwidth.py + spacedlines.py + spaced.tsv, zoom.py, pagecheck.py + PAGECHECK-BRIEF.md,
%  patches.py (PATCHES, NOTES, POST), decisions.tsv, the sweeps, split5.py, rebuild.sh.
% =====================================================================
\documentclass[12pt,a4paper,openany]{book}
\usepackage[utf8]{inputenc}
\usepackage[T1]{fontenc}
\usepackage{libertinus}
\usepackage{libertinust1math}
\usepackage[danish]{babel}
\usepackage[protrusion=true,expansion=true]{microtype}
\usepackage{setspace}
\usepackage[margin=1.3in]{geometry}
\usepackage{fancyhdr}
\usepackage{hyperref}
\hypersetup{pdftitle={Om Selvforstaaelse i det Onde}, pdfauthor={Rasmus Nielsen}, hidelinks}
\pagestyle{fancy}
\fancyhf{}
\fancyhead[LE,RO]{\thepage}
\fancyhead[C]{\textit{R. Nielsen: Om Selvforstaaelse i det Onde (1869)}}
\setstretch{1.3}
\usepackage{marginnote}
\renewcommand{\thefootnote}{\arabic{footnote})}   % the print numbers each page's note ¹) (p. 553: *))
\newcommand{\opage}[1]{\marginnote{\footnotesize\textit{[#1]}}}
% the essay's heading: title (and subtitle), centred, a short rule; the print's large initial is not reproduced
\newcommand{\essayhead}[2]{\begin{center}{\large #1}\ifx&#2&\else\\[0.3em]{\normalsize #2}\fi\\[0.5em]\rule{3em}{0.4pt}\end{center}\medskip\noindent}
\newcommand{\signature}[1]{\par\smallskip\hfill #1\hspace*{2em}\par}

\begin{document}

% --- front matter: the volume's title page (PDF 7), read at the image (reader K, F7; the page check reads it again blind).
\begin{titlepage}
\begin{center}
\vspace*{3em}
{\Huge For Ide og Virkelighed.}\\[1em]
\rule{8em}{0.4pt}\\[1em]
{\large Et Tidsskrift,}\\[0.8em]
{\small udgivet}\\[0.4em]
{\footnotesize af}\\[0.6em]
{\large R. Nielsen,}\\[0.4em]
{\large Bj. Bjørnson og Rud. Schmidt.}\\[1em]
\rule{8em}{0.4pt}\\[1em]
{\large Aargang 1869.}\\[1em]
\rule{8em}{0.4pt}\\[4em]
Kjøbenhavn.\\[0.4em]
Chr. Steen \& Søns Forlag.\\[1.5em]
{\small Louis Kleins Bogtrykkeri.}
\end{center}
\end{titlepage}

% ===== TEXT (pp. 521-554) =====
\essayhead{Om Selvforstaaelse i det Onde,}{med Henblik paa Shakspeares Richard den Tredie.}
% --- p. 521 ---
\setcounter{footnote}{0}%
\opage{521}Naar en Rovfugl slaaer ned paa et Sødyr, og de stridende
Kræfter paa begge Sider saa omtrent ere lige,
vil Kampens Udfald fornemmelig afhænge af det Element,
hvori der kæmpes. Lykkes det Havets Indvaaner
at drage sin vingede Fjende ned under Bølgerne, da
vinder den en dobbelt Fordel, den strider i sit Element,
og Elementet strider for den. Det er nu ikke blot i
Dyrelivet, at Kræfternes Styrke i Striden betinges af
Elementet; ogsaa i Aandslivet viser sig et Tilsvarende.
De aandige Kræfter i Tro og Viden ere med Hensyn
paa Principer og Forudsætninger, Forestillinger og Begreber
saa forskjellige, at de tilsvarende Ideer maae
siges at bevæge sig i to absolut uensartede Regioner,
og Striden imellem dem hvergang beroe paa det Element,
hvori den føres. At drage de religiøse Livstanker
ned i Videns Strøm, er vilkaarlig at tage Parti for det
Antireligiøse; thi Viden strider i sit eget objektive Element,
og Elementet strider for den. Det er et ældgammelt
Misgreb, et tusindaarigt theologisk Bedrag, at forvexle
Elementerne, forgude det Objektive og drukne
Troens Ørn i Videnskabens Vande.

% --- p. 522 ---
\setcounter{footnote}{0}%
\opage{522}Det er en Grundsætning, som Enhver, der vil, kan
forstaae, at en Samvittighedssag er religiøs og sædelig,
men ikke videnskabelig; eller er der maaske noget fornuftigt
Menneske, som i det virkelige Liv forvexler et
Samvittigheds-Anliggende med et videnskabeligt Problem?
Dog vel heller ikke omvendt! Naar en Jesper Ridefoged
føler i sin Samvittighed, at Jorden ikke er rund,
saa leer man; og hvorfor? Fordi man uvilkaarlig føler,
at Fysik og Samvittighed maa tale hver sit Sprog, have
hver sin Tankekreds med egen Begrundelse af egne
Grunde. At videnskabelige Spørgsmaal ikke umiddelbart
kunne underkastes Samvittighedens Kjendelse,
betyder, at Videnskaben selv er upersonlig, at Videns
Element er objektivt. De religiøse og ethiske Sandheder
ere saavel i sidste som i første Instans Samvittighedssandheder;
al sand og grundig Forstaaelse af
disse Sandheder beroer paa Selvforstaaelse, men Selvforstaaelsens
Element er Subjektiviteten. Hermed er det
Objektive i Religion og Sædelighed ingenlunde svækket,
endnu mindre fornægtet; kun dette skal fremhæves, at
Dommen derover er subjektiv, thi Dommeren er i Menneskets
Selv. Overhører man denne indre Dommer eller
nedsætter ham til en underordnet Betjent og søger
Overdommeren i det Videnskabelige, da er det let at
forudsee, hvad Følgen maa blive. Det er saaledes ikke alene
for Religiøsitetens, men ogsaa for Kulturhistoriens og Humanitetens
Skyld af stor Betydning at gjøre sig Rede
for Forskjellen mellem Videnskabssandhed og Samvittighedssandhed.

I alle Lande, hvor Videnskab og Literatur holde
Skridt med Tiden, foreligge nu literære Aktstykker, som
% --- p. 523 ---
\setcounter{footnote}{0}%
\opage{523}klarlig bevise, at Theologien med sin lærde Oppakning
er uden Princip, at dens middelalderlige Rolle er udspillet.
Men, vil man maaske indvende, dersom dette
forholder sig saaledes, hvorledes skal man da forklare
sig, at en saa vitterlig, for Kultur og Humanitet saa
dybt indgribende Kjendsgjerning endnu ikke ret er gaaet
op for den dannede Almenhed? Sagen er let at forklare.
Ved at tvinge Troen ned i Videns objektive Element,
har Theologien fra Slægt til Slægt hexet sig selv og Samtiden
ind i den Fixide, at Theologi og Religion ere uadskillelige;
»falder Theologien, saa falder Religionen«, det er det
dogmatiske Løsen, der skal indjage den civiliserede Verden
Skræk for Atheisme og Respekt for Theologi. Hvor
falsk dette Løsen er, kan sees af Striden om Tro og
Viden. Theologiens og Atheismens øverste Grundsætning
er nemlig een og den samme. »Falder Theologien,
saa falder Religionen«, kun af denne Forudsætning lader
det sig forklare, at en atheistisk Kritik, der videnskabelig
har opløst Theologien, kan bilde sig ind, at den med
det Samme har opløst den positive Religion og saaledes
vundet en dobbelt Seir. At Theologiens »objektive Tro« er
falden for Videns Vaaben, er vitterligt for alle Vidende, og
dog, mærkværdigt nok, en offenlig Hemmelighed. Hvorledes
en saadan Hemmelighed kan holde sig i en oplyst
Tidsalder, trods saamange Kikkerter og saamange klare
Øine, kan vel forstaaes, naar man betænker, at det ikke
blot er i Theologernes, men ogsaa i Atheisternes velforstaaede
Interesse, at det menige Folk, »den store
Hob«, ikke kommer til at kige for dybt i Gryden. Fra
begge Sider finder man sin Regning ved at dække Kjendsgjerningen
med et tyndt Flor. Hvad Theologien angaaer,
% --- p. 524 ---
\setcounter{footnote}{0}%
\opage{524}vilde det jo være et Selvmord, om den udbasunede sit
eget Nederlag og lod ved Avertissementer bekjendtgjøre,
at den virkelig er en forloren Videnskab; det bliver
Modstandernes Sag at sørge for Bekjendtgjørelsen. Man
skulde nu mene, at blandt disse Modstandere maatte
Atheisterne med den dobbelte Seir staae i forreste
Række. Saaledes var det da ogsaa i forrige Aarhundrede;
men i det nærværende har Situationen mærkelig
forandret sig. Man har gjort en psykologisk Erfaring.
Atheisten vil, naar han ret besinder sig, indsee, at han
ikke kan være tjent med, at den theologiske Troesvidenskab
i Folks Omdømme taber al Kredit, thi naar den
Fjende, han lever af at bekæmpe, er død og borte, saa
har han væsenlig Intet mere at gjøre; hans esoteriske
Viden er en negativ Magt, der ikke tør blotte sit positive
Grundlag. Saalænge Atheismen har det metafysiske
Slør over sig, er den ikke ulig Gudinden i Sais;
man føler en Gysen for den mørke Hemmelighed, den
gaadefulde, forstenende Sandhed. Men naar Sløret løftes,
er det forbi med Ligheden. Ynglingen bliver ikke
sindssvag af Forfærdelse, tvertimod, han bliver vranten,
fordi der ingen Forfærdelse er, fordi hans Forventning
skuffes: »Er dette Sandheden selv, er det Gudindens
Aasyn? Et saadant Ansigt kan man jo træffe paa enhver
Dandsebod«! Den platte Fornægtelse kan nøies
med Vidnesbyrd af de fem Sandser. For at forstaae
den Lære, at der ingen personlig Gud er og ingen
Sjælens Udødelighed, behøver den praktiske Kulturbevidsthed
ikke at plage sig med dunkle filosofiske Talemaader;
en saadan Almindelighedens Visdom indaander
man med den almindelige Luft. Af Atheismen tør man
% --- p. 525 ---
\setcounter{footnote}{0}%
\opage{525}altsaa ligesaalidt vente en fuldstændig Afsløring som af
Theologien. Under saadanne Omstændigheder maa den
offenlige Mening, hvad enten den nu til en given Tid
helder til Theologien eller til Atheismen eller maaske --- eftersom Udstillingen er i »Butikken« --- til begge
paa een Gang, føle med sig selv, at den er uden Indsigt,
uden Overbevisning, kun en Mening, og saa kan Hemmeligheden
nok blive bevaret.

Den offenlige Hemmelighed har imidlertid et Gode
ved sig, som vi ikke bør miskjende. Der gaaer igjennem
den civiliserede Verden en Fornemmelse af, at Striden
for og imod de religiøse Dogmer er upraktisk, at
det Praktiske i Sagen, det, hvorpaa det virkelig kommer
an, er Sædelighed og sædelig Selvforstaaelse; thi,
hedder det, »den, hvis Vandel er ren og ret, hans Tro
kan umulig være slet«. Theorier, der angribe den aabenbarede
Religion --- de kalde sig nu Deisme, Pantheisme,
Atheisme eller, hvad de helst ville --- ere ingenlunde
farlige for Samfundet, naar de kun ikke angribe Sædeligheden.
Fra Statens og Humanitetens Synspunkt er
en dygtig, sund og kraftig Sædelære langt at foretrække
for dunkle, mystiske Aabenbaringsspekulationer. Lad,
siger man, en sindig filosofisk Tænker levere os en forsvarlig
Ethik, en praktisk Sædelære, der hævder Pligt og
Ret og gaaer tilbunds i Modsætningen af Godt og Ondt,
og han skal være os velkommen: vi skulle visselig ikke
genere ham med nærgaaende Spørgsmaal om Troesbekjendelse.
At den Anskuelse, man i populær Forstand saa
let kan stemple som Atheisme, ikke just, hvad det Sædelige
angaaer, altid staaer tilbage for orthodoxe Lærebegreber,
har Exemplet med den ældre Fichte tilstræk%
% --- p. 526 ---
\setcounter{footnote}{0}%
\opage{526}kelig vist. Denne store, sædelige, folkelig sindede Tænker
fik man, da man ikke paa anden Maade kunde komme
ham tillivs, omsider anklaget for Regeringen som Atheist.
Det hjalp, at sige i borgerlig Forstand, thi han blev
som Professor i Filosofien ved Universitetet i Jena virkelig
afsat fra sit Embede; men hvad det betød i videnskabelig
og sædelig Forstand, kan sees af hans mandige,
aabne Forsvar. Fordi et Menneske opgiver Troen
paa en personlig Gud, derfor opgiver han ikke Troen
paa en guddommelig Verdensorden. »Denne Tro«, siger
Fichte, »maa man ikke forestille sig som en vilkaarlig
Antagelse, som Mennesket kan gjøre eller ikke, eftersom
det behager ham, som en fri Beslutning at holde for
sandt, hvad Hjertet ønsker, fordi det ønsker det, som
en Fuldstændiggjørelse eller Erstatning af overbevisende
Grunde ved at tage Haabet til Hjælp. Hvad der er grundet
i Fornuften, er slet og ret nødvendigt; og hvad der ikke
er nødvendigt, er netop af den Grund fornuftstridigt. At
antage Sligt er Indbildning og Drøm, hvor fromt man
endogsaa drømmer. Det sande Gode erkjendes kun under
den Betingelse, at jeg finder mig fri for alle Sandseverdenens
Indflydelser, absolut virksom i mig selv,
altsaa som en over alt Sandseligt ophøiet Magt. Denne
Frihed er ikke ubestemt; den har sit Maal, kun at den
ikke erholder det udenfra, men sætter sig det ved sig
selv. Paa denne Frihed og denne dens Bestemmelse
kan jeg ikke tvivle uden at opgive mig selv. Jeg kan
ikke tvivle, siger jeg --- kan ei engang tænke mig Muligheden
af, at det ikke skulde være saaledes, at hin
indre Stemme kunde bedrage, at den først maatte autoriseres
og begrundes andenstedsfra. Her gjælder ingen
Ræsonneren, ingen Omforklaren eller konstlet Udtydning.
% --- p. 527 ---
\setcounter{footnote}{0}%
\opage{527}Hin Stemme er absolut positiv og kategorisk. Jeg kan ikke
komme videre, naar jeg ikke vil forstyrre mit eget Indre;
jeg kan just derfor ikke gaae videre, fordi jeg ikke kan
ville gaae videre. Her ligger det, som sætter Ræsonnementets
ellers utæmmede Flugt sin Grændse, det, som
binder Aanden, som binder Hjertet; her er det Punkt,
som forener Tænken og Villen til Eet og bringer Harmoni
i mit Væsen. Dette er den sande Tro; den moralske
Orden er det \emph{Guddommelige}, som vi antage.
Det konstrueres ved at handle ret. Dette er den eneste
mulige Troesbekjendelse: freidig og uhildet at fuldbringe,
hvad Pligten hvergang byder, uden Tvivlen og Ræsonneren
over Følgerne. Derved bliver det Guddommelige
os levende og virkeligt, enhver af vore Handlinger fuldføres
i dettes Forudsætning, og Følgerne opbevares deri«.

Man seer af denne Udtalelse, hvor Fichte vil hen,
og hvad det er, hvori hans sædelige Bevidsthed grunder.
Dersom Fichte nu blot var enig med sig selv om det
Element, hvori det Godes Ide skal røre sig og leve, saa
kunde vi trygt følge hans Veiledning; thi at han for
Alvor vil hævde Troen paa det Gode som det i sig
Sande og Guddommelige, derom er ingen Tvivl. Men
netop i Sammenblandingen af Troens og Videns Sfære,
altsaa i Elementernes Forvexling ligger hans Selvmodsigelse.
Han forkaster Troen paa en personlig Gud;
det er ingenlunde af Trods, Vantro eller Ugudelighed,
det er i Kraft af videnskabelige Grunde, det er fordi en
personlig Gud, ifølge hans Begreb, maa være en indskrænket
Gud og en indskrænket Gud utænkelig. Vi
høre ham i sit Forsvar: »I tillægge Gud Personlighed
og Bevidsthed. Hvad forstaae I da ved Personlighed og
Bevidsthed? Dog vel det, som I have fundet i Eder
% --- p. 528 ---
\setcounter{footnote}{0}%
\opage{528}selv og betegnet med dette Navn? Men at I paa ingen
Maade tænke eller kunne tænke Eder dette uden Begrændsning
og Endelighed, det kan den ringeste Opmærksomhed
ved Eders Konstruktion af dette Begreb
lære Eder. Ved at tillægge dette Væsen hint Prædikat
gjøre I det altsaa til et endeligt, til et Væsen, der er
Eder ligt, og I have ikke, som I vilde, tænkt Gud, men
kun mangfoldiggjort Eder selv i Tænkningen«. Men
spørge vi nu vor Filosof, af hvad Grund vi da skulle
troe paa det Godes Realitet og paa en guddommelig
Verdensstyrelse, saa synes det, som om han med Eet
vilde flytte os over i en helt anden Sfære, et andet
Element. »Overbevisningen om vor moralske Bestemmelse«,
siger han, »gaaer allerede selv frem af moralsk
Stemning, og er Tro, og man siger forsaavidt ganske
rigtigt: Elementet for al Vished er Tro. --- Saaledes
maatte det være; thi Moralitet, saa vis den ogsaa er,
kan simpelthen kun konstrueres ved sig selv, ingenlunde
ved nogen logisk Tanketvang«. Her seer man Inkonsekvensen:
Spørgsmaalet om Guds Personlighed skal afgjøres
i Videnskaben, herover tør man altsaa nok ræsonnere;
men Spørgsmaalet om Moralitetens Gyldighed
skal ikke afgjøres i Videnskaben, thi derover tør man
ikke ræsonnere. Inkonsekvensen bliver endnu større, naar
vi lægge Mærke til, at Forfatteren efter at have fastsat
Selvets Princip subjektivt i Troen --- altsaa udenfor
Videnskaben --- ikke destomindre vil have dets Udvikling
af »det slet og ret Nødvendige« ført over i Videnskaben.
Naar altsaa En siger: Jeg kan ikke opgive
Moraliteten uden at opgive mig selv, da erklæres han
for Videnskabens Tiltale fri at være; her gjælder ingen
logisk Tanketvang; men dersom Nogen siger: Jeg kan
% --- p. 529 ---
\setcounter{footnote}{0}%
\opage{529}ikke opgive den personlige Gud uden at opgive mig selv,
da bør Videnskaben stævne ham, thi den personlige Gud
er afsat ved »logisk Tanketvang«.

Kun sædelig Iver for at gjennemføre en ubetinget
rationel Selvforstaaelse i det Gode har kunnet forlede
en Tænker som Fichte til at begaae en saa iøinefaldende
Inkonsekvens. Det ethiske Misgreb, som her foreligger,
kan bedst belyses, naar vi, istedenfor at tale almindeligt
om det almindelige Gode og overveie, hvorvidt der
gives en rationel Selvforstaaelse i det Gode, kaste et Blik paa
det Godes Modsætning og spørge, om der da ikke ogsaa gives
en rationel Selvforstaaelse i det Onde. Saalænge Talen
blot er om det Gode, slaaer man sig let til Ro ved almindelige
Betragtninger. Det Gode, hedder det, er Eet
med Menneskets sande Natur, idet jeg vælger det Gode,
vælger jeg væsenlig mit eget Selv; at jeg gjør det Gode
for det Godes Skyld, betyder, at jeg i dybeste Forstand
gjør det for min egen Skyld osv. Da det nu gaaer saa
glat med slige skjønne, almindelige Talemaader, saa bliver
man let fristet til at behandle Samvittighedssandheder
som videnskabelige Læresætninger og drage Personlighedens
høieste Anliggender ned i et upersonligt
Element. Anderledes, naar Talen er om det Onde. Her
ville de rene Almenbetragtninger ikke slaae til. Skal
det Ondes Væsen begribes i Videnskaben, da løber Forklaringen
tilsidst ud paa, at det Onde kun er et Bedrag,
et Skin. At forstaae sig selv i det Onde, maatte
da, ifølge denne Synsmaade, ansees for ligesaa umuligt
som at gjøre det Onde, fordi det er ondt. Kun ved at
ombytte de almene Afstandsbetragtninger med et virkeligt
Nærsyn, det vil sige: med en til Ide og Virkelighed
svarende Opfattelse af det Personlige i Karakter
% --- p. 530 ---
\setcounter{footnote}{0}%
\opage{530}og Villie, vil det være Tanken muligt at naae til Roden
af det Onde. Et fra dette Synspunkt oplysende Karakterbillede
have vi i Shakspeares Richard den Tredie,
Ridderen af det radikale Onde, en kongelig Satan med
Kjød og Blod, en inkarneret Djævel.

Det Egne ved Richards Karakter er ikke det, at
han er blodtørstig, grusom og begaaer Forbrydelser; i
denne Henseende er han kun en Søn af sin Tid. Den
Tid, hvori Richard lever, er en ond Tid: Vold og Frækhed,
Mord og Mened høre til Dagens Orden. Den middelalderlige
Stat med sit stærke, selvraadige Aristokrati
er i det femtende Aarhundrede given til Pris for Partilidenskaber,
der true den med Undergang; Striden mellem
den hvide og den røde Rose, en mere end trediveaarig
Borgerkrig, i hvilken tredsindstyve Personer af
Kongehuset og Halvdelen af den engelske Adel tilsatte
Livet, er et blodigt Afsnit i Englands Historie. Nei, at
Richard er en Tyran, er ikke det Karakteristiske; det
Karakteristiske ligger i den Bevidsthed, hvormed han
udøver sit Tyranni. De handlende Personer i hans Omgivelse,
hvilke Ugjerninger de end begaae, have dog
et Slags Undskyldning i deres ophidsede Lidenskaber;
indhvirvlede i Øieblikkets Hævnkrig, ere Alle mere eller
mindre forblindede, kun ikke Richard. Det er et Træk
i hans forfærdelige Overlegenhed, at medens Ingen iblandt
de Medstridende, Ven eller Fjende, ret er i Stand til at gjennemskue
ham, gjennemskuer han dem alle, seer ned paa
dem, foragter dem og bruger dem som blinde Redskaber;
ja, hvad mere er --- thi hans dæmoniske Blik er
ligesaa indadvendt som udadvendt --- idet han gjennemskuer
de Andre gjennemskuer han tillige sig selv
% --- p. 531 ---
\setcounter{footnote}{0}%
\opage{531}og har, som det synes, en fuldstændig Selvforstaaelse i
det Onde.

Monologen, hvormed første Akts første Scene begynder,
lader os gjøre et dybt Indblik i Richards Sjæl.
Krigens Tummel er hans Element, han kjeder sig i Freden;
Broderen Edvard den Fjerdes blødagtige Hof er
ham modbydeligt. Ikke af sædelig Indignation, men fordi
han har sine egne Planer, og saa fordi han, »hæslig
myntet«, stedmoderlig behandlet af Naturen, ingen Lykke
gjør hos Hoffets Damer, erklærer han, at han vil være
Nidding »og hade disse Dages tomme Skæmt«. --- »Anslag
og Sammenrottelser jeg danned' ved mørke Varsler,
Drømme og Libeller, som skal min Broder Clarence og
Kong Edvard til Had og Fjendskab mod hinanden ægge,
og om Kong Edvard holde vil sit Ord, saavist som jeg
er listig, falsk og troløs, skal Clarence end idag i Fængsel
spærres.« Saa høre vi strax fra Begyndelsen, hvad
Richard er, og hvad han vil: han er en besluttet Mand,
en Mand, der er enig med sig selv. Det gruopvækkende
Forsæt, hvormed Richard optræder, og den Maade,
hvorpaa han motiverer det, har foranlediget en Misforstaaelse,
som vi her maae afvise. Man har forstaaet hans
Ord, som om han i sin naturlige Hæslighed søgte en
Undskyldning for den Ondskab, hvormed han vil hævne
sig paa det menneskelige Samfund, fra hvilket han som
Vanskabning er udelukket. Saaledes A. W. Schlegel.
»Ondskab«, siger denne Forfatter, »er ikke Andet end
en med Forsæt samvittighedsløs Egoisme; men alligevel
kan den aldrig ganske aflægge Sædelighedens Form, fordi
denne er Lov for alle tænkende Væsener; den maa søge
at henføre sin vrange Handlemaade til Grundsætninger.
% --- p. 532 ---
\setcounter{footnote}{0}%
\opage{532}Uagtet Richard meget vel erkjender sit Sindelags Afskyelighed
og sin Sendelse fra Helvede, søger han dog
at retfærdiggjøre sig for sig selv ved Skingrunde; da
den Lykke at elskes er nægtet ham, hvad bliver ham
saa tilbage uden den Lykke at herske?« Denne Bemærkning,
der forøvrigt ikke er uden en vis psychologisk
Finhed, passer ikke paa Richard; den passer paa
forsmaaet Kjærlighed, men den, der aldrig har elsket
Nogen uden sig selv alene, kan ikke føle en forsmaaet
Kjærligheds Kvaler. Den Kjærlighed, hvorom Richard
taler, er da heller ikke af den Natur, at den kan betragtes
som en Dyd. Naar han siger, at han ikke kan »spille Elsker«,
at han »ei kan beile til sig selv i Speilet« og
»ei er skabt til strunk at kneise for en Nymfe kaad«,
da mener han ingenlunde, at det er noget Dydigt
eller Godt, der gaaer tabt for ham. Richard tænker
overhovedet slet ikke paa en Skælnen mellem Dyd
og Last, men snarere paa en Vælgen mellem forskjellige
Laster; han er ikke skikket til den ene Last, og
derfor vælger han den anden; --- »da jeg Elsker ei kan
spille og underholde disse vakkre Damer, har jeg besluttet,
jeg vil være Nidding;« det er i Form af Undskyldning,
som om En vilde sige: »Da jeg ei kan stjæle,
har jeg besluttet, at jeg vil være Røver.« Det er bogstaveligt,
at Richard veed, hvad en Nidding er, og \emph{desuagtet}
eller rettere \emph{desaarsag} har besluttet at være
Nidding. »Saa vil han altsaa det Onde, fordi det er
ondt, hvilken uhyre Modsigelse!« Ganske vist er det
en ubegribelig Modsigelse, saalænge vi betragte ham fra
Videns almene objektive Standpunkt, men ingenlunde,
naar vi forstaae ham personligt. Saalænge Ordet Nid%
% --- p. 533 ---
\setcounter{footnote}{0}%
\opage{533}ding er en Betegnelse for den Onde i streng Modsætning
til den Gode, med andre Ord: saalænge Forskjellen
mellem Godt og Ondt endnu ikke er udslettet fra
Grunden af, saalænge vil det ganske vist være umuligt
for et Menneske at ville det Onde, fordi det er ondt.
Men Forskjellen mellem Godt og Ondt har Richard
personlig udslettet. Ved at udslette denne Forskjel er
han ingenlunde kommen til at staae paa neutral Grund
som En, der hverken er god eller ond; han er tvertimod
bleven fordybet i det Onde; thi at nægte Forskjellen
mellem Godt og Ondt, er at vie sig til det Onde. Hvorledes
er da Richard oprindelig kommen til at hæve
Skilsmissen mellem Godt og Ondt? Ved at fornægte
Gud, altsaa ved at stille sig paa et atheistisk Standpunkt! Men med Richards Atheisme har det en egen Sammenhæng;
den skyer det objektive Element og holder
sig konsekvent i det subjektive. Richards Grunde ere
ikke videnskabelige, men ubetinget personlige. Til at
skrive en kritisk metafysisk Piece, hvori han afskaffede
den aabenbarede Religion og bebudede en ny Æra,
da Menneskeslægten skulde lyksaliggjøres med »en immanent
Udvikling af sit immanente Princip«, havde han jo
ingen Tid; at indføre en ny Liturgi med Formular: der
er ingen Gud, og Richard er hans Profet, ansaae han
ikke for hensigtsvarende, saa tog han Sagen subjektivt,
saa subjektivt, at han nøiedes med at »blive sig selv
gjennemsigtig«, uden at blive gjennemsigtig for Andre.
Den indre Vished, hvormed Richard har opgivet Gud,
fremlyser netop af den Frækhed, hvormed han hykler
Gudsfrygt og bruger Religionen, hvor den kan tjene
ham som Middel, til sit Øiemeds Opnaaelse. Hans Tan%
% --- p. 534 ---
\setcounter{footnote}{0}%
\opage{534}ker desangaaende udtales tydelig nok i Monologen, (Første
Akt, tredie Scene mod Slutningen): »Jeg Skaden
gjør, jeg først begynder Striden; det hemmelige Onde
jeg udbreder, dets Byrde lægger jeg paa Andres Skuldre\dots{}.
Nu troe de det, og stedse hvæsser mig til
Hævn mod Rivers, Vaughan og mod Grey. Da sukker
jeg og viser med et Skriftsprog, at Gud os byder Ondt
med Godt at lønne. Saaledes skjuler jeg den nøgne
Nidding med gamle Stumper stjaalne ud af Biblen, og
synes Helgen, naar jeg spiller Djævel«. I hvilken Grad,
hans Hykleri kan brillere, see vi da især af den guddommelige
Komedie, han spiller med Londons Borgere.
(Tredie Akt, syvende Scene.) Clarence er myrdet; Edvard
er død, og hans unge Sønner forvarede i Tower,
Rivers, Grey og Vaughan og selv den ædle, umistænksomme
Hastings have alle besteget Skafottet; saaledes
har Richard forberedt Alt, og nu gjælder det kun om,
at hans Thronbestigelse kan faae Skinnet af at skee paa
Folkets Opfordring og være saa legitim som muligt.
Richard maa naturligviis lade, som om han saare nødig
tog Regeringens tunge Byrde paa sig. Greben af mørke
Anelser ville Borgerne da ogsaa nødig til at vælge ham.
Men Buckingham, Richards Kopi og Redskab, veed at
bygge Fremtid paa »en hellig Plan«. »Lad«, siger han til
Richard, »som I frygtsom er, lad Eder længe bede, før
I sees. Stil Eder, Mylord, mellem tvende Klerke og tag en
Bønnebog i Eders Haand«\dots{} Da nu Borgerne med deres
Lord Mayor og Aldermænd komme, blive de naturligvis
afviste af den fromme, i gudelige Betragtninger fordybede
Richard. For at de gode trofaste Sjæle dog ikke
ganske skulle tabe Modet, fordi den »ædle Hertug«
% --- p. 535 ---
\setcounter{footnote}{0}%
\opage{535}nøler med at modtage dem, forklarer Buckingham nu Sagens
egenlige Sammenhæng, idet han siger: »Nei, Mylord,
denne Prinds er ingen Edvard; han sig ei strækker
paa den bløde Løibænk, men paa sit Knæ i gudelig Betragtning.
Han lefler ei med et Par kaade Friller, men
med to Geistlige han gransker dybt.« Da saa Lord
Mayoren omsider seer Richard paa en Altan mellem to
Biskopper, forsvinder al Tvivl, og han udbryder i from
Begeistring: »See hist hans Naade mellem tvende Bisper,
to Dydens Støtter for en kristen Prinds, som
værne mod Forfængeligheds Glimmer! og see en Bønnebog
er i hans Haand, den sande Pryd, hvorpaa man
Fromhed kjender« \dots{} For at blive, hvad han siden selv
kalder sig, en »Herrens Salvede«, kan Richard ret føle
Trang til at granske med to Geistlige, og, om saa skal
være, gaae til Skrifte hos to Bisper: en Mand, der daglig
skrifter saa oprigtig for sig selv, kan vel have et Par
Skriftemaal tilovers for Bisperne. Han er en fuldendt
Hykler; Hykleri er en Art Gudsfornægtelse, det fuldendte
Hykleri en gjennemført Gudsfornægtelse.

Altsaa: Richard er Atheist; fra dette Standpunkt
maa hans Selvforstaaelse og Konsekvenserne i hans Karakter
bedømmes. Viser det sig nu, at en atheistisk
Ethik ikke er i Stand til at magte Richard, da er det jo
med det Samme bevist, at den atheistiske Ethik ikke
eier det Princip, hvoraf en sædelig Modsætning mellem
det Gode og det Onde lader sig udvikle og hævde. Vi
skulle nu prøve, hvorledes Sagen forholder sig.

Naar Gud er fornægtet, altsaa Gudsvillien udelukket,
saa er Mennesket, seet i Ideen, ubetinget sin egen
Herre, sin egen Lovgiver, sin egen Dommer. Dette
% --- p. 536 ---
\setcounter{footnote}{0}%
\opage{536}maa naturligvis ikke forstaaes, som om Individet i sit
Indvortes --- thi hvad det Udvortes angaaer, behøves
her ingen Forklaring --- kunde vilkaarlig gjøre sig selv til
hvad det vil: den individuelle Villie er betinget af og forsaavidt
bunden til den individuelle Natur, og Karakterens
Udvikling paa een Gang en personlig Udvikling baade af
Natur og Villie. Det nytter i denne Sammenhæng ikke at
skyde sig ind under det metafysiske Almene og tale vidt
og bredt om det Almenmenneskelige, den almene Natur,
den almene Fornuft, den almene Tanke, den almene Villie;
thi i ethvert virkeligt Menneske er det Almene begrændset
og antager eiendommelig Skikkelse; derfor har
enhver menneskelig Individualitet jo sin Natur, sin Fornuft,
sin Tanke og sin Villie: det Almene er ikke udelukket,
men det er individualiseret. Vi kunne nu ikke
sige, at den virkelige Menneskenatur, Naturen existerende
i virkelige Mennesker, er enten ubetinget god
eller ubetinget ond; den er, som bekjendt, en Blanding
af begge. Som vi maae antage, at der i den onde Natur,
hvor mørk og fordærvet den end er, dog altid maa
være noget Godt, efterdi en absolut ond Natur vilde
være Unatur, og en absolut Unatur en Umulighed, saaledes
maae vi jo ogsaa antage, at der i den gode Natur,
hvor ædel og ren den forøvrigt kan være, saasandt
den er af adamisk Princip, maa være noget Ondt.
Men saa er hermed udtrykkelig indrømmet, at det Onde
i Menneskenaturen ligesaavel har sin Almindelighed, som
det Gode. Men naar det Almindelige ligesaavel kan
svare til det Onde som til det Gode, saa er det jo forgjæves
at ville hævde den sædelige Fordring ved at henvise
Individet til det i Menneskenaturen Almindelige.
% --- p. 537 ---
\setcounter{footnote}{0}%
\opage{537}Forskjellen mellem den gode og den onde Natur er relativ
og har til sine Extremer paa den ene Side saadanne
Individer, i hvilke det Gode har oprindelig Overvægt
over det Onde, paa den anden Side saadanne, i
hvilke det Omvendte er Tilfældet; denne Forskjel faaer en
gjennemgribende Betydning, naar Opgaven er at bestemme
det sædelige Princip. Høre vi igjen Fichte, den Mand
om hvem Goethe, skjøndt hans fine Verdensklogskab ikke
kunde billige Ethikerens energiske Optræden, selv tilstaaer,
at han var »en af de dygtigste Personligheder,
man nogensinde har seet«, en Mand, paa hvis »Tænkemaade
der i høiere Henseende Intet er at udsætte«: saa
føle vi, at Fichte taler ud af sin Sjæls Inderste, naar
han erklærer, at han ikke kan tvivle paa den sædelige
Frihed og dens Bestemmelse, uden at opgive sig selv,
og saaledes peger paa »det Punkt, som forener Tanken
og Villien til Eet og bringer Harmoni i hans Væsen.«
Han vil det Gode, som han vil sin egen sande Villie; han
kan ikke gaae videre, naar han ikke vil forstyrre sit eget
Indre; han har sin Selvforstaaelse i det Gode, fordi
hans Natur er god. Vende vi os dernæst til Richard,
da faae vi samme Forklaring med samme Begrundelse
og dog i aldeles modsat Retning. Ogsaa Richard har
Loven for sine Handlinger ubetinget i sig selv; ogsaa
han vil, hvad han ifølge sine eiendommelige Naturanlæg
maa ville. Saa lidt som en Fichte kan forandre
sin Natur, ligesaa lidt kan en Richard forandre sin;
forsøger han derpaa, saa synker han og gaaer øieblikkelig
tilgrunde. Fra dette Synspunkt faaer den Selvbetragtning,
hvori han med øiensynligt Velbehag fremhæver
sin Skabnings Hæslighed --- »Af Mandens Billed
% --- p. 538 ---
\setcounter{footnote}{0}%
\opage{538}er jeg en Forkortning; bedragerisk den hyklende Natur
beskrev mit Ansigtstræk. Jeg ud i Verden fortidlig
sparket blev af Moders Liv, halv færdig knap, saa lam,
saa ufuldendt, at Hunde bjeffe, naar jeg hos dem tøver« ---
en anden Betydning end den, at tjene til tom Undskyldning,
fordi han ikke vil, hvad Menneskene kalde det
Gode. Richard beklager sig ikke, undskylder sig ikke;
han tager sig selv, som han er, han ønsker ikke at
være anderledes eller at være en Anden: han er »sig
selv alene«. Hans ydre Hæslighed er Symbol paa et
djævelsk Indre; men han er ingenlunde misfornøiet med sit
Ydre; han har ingen Grund dertil, tvertimod! Vanskabningen
er for den dæmoniske Richard, hvad det frastødende
Ydre, Braknæsen, de fremstaaende Øine, Vorterne,
den hængende Mave og de plumpe Fødder, var for den
ironiske Sokrates, --- en ypperlig Maske, der passer til
hans Rolle. Spørgsmaalet er her ikke om, hvad der er
hæsligt og hvad der er skjønt, saalidt som om, hvad der er
Ret og hvad der er Uret: Magt er Ret. Richards Ret
er Egenvilliens Magt i Sindets Kræfter, og der er Kræfter
i Richard. Aanden vil frem, den mørke Aand, der
gløder i hans Herskesyge, hans bittre Følelse af Had
og Hævn; der brænder i ham en Vulkan, og Vulkanen
maa komme til Udbrud. Atheistisk kan altsaa Richard
i denne Henseende siges at forstaae sig selv i det Onde
ligesaa ubetinget som Ethikeren forstaaer sig selv i det
Gode: Retningerne ere modsatte, men Standpunktet er
det samme, Principet det samme.

For at hævde den principielle Forskjel mellem Godt
og Ondt maa den atheistiske Ethik nu søge at give Sagen
en anden Vending. Den begyndte med at opgive
% --- p. 539 ---
\setcounter{footnote}{0}%
\opage{539}»den logiske Tanketvang«, denne Overilelse maa den
bitterlig fortryde. Skal den kunne magte Richard, maa
den begynde, ikke personligt med subjektiv Enhed af
Fornuft og Natur, men upersonligt med objektiv Enhed
af Fornuft og Tænkning. Er det sandt, hvad Fichte
siger, at den sædelige Selvbevidsthed finder sig »fri for
alle Sandseverdenens Indflydelser, absolut virksom i sig
selv, altsaa som en over alt Sandseligt ophøiet Magt«,
da maa den opgive Existensen, tye ind under Metafysikens
»Almene« og føre Striden over i Videns logiske
Element. Det Selv, hvorpaa Ethikeren tænker, naar han
i det Godes Navn taler om at kjende sig selv og ville
sig selv af ren Nødvendighed, er da ikke det individuelle,
men det almene. Kun det sandselige, driftbundne Selv
er individuelt; det sande aandelige Selv er universelt,
det er en i sig gjennemsigtig Theori. Ikke hvad Individet
vil, men hvad Fornuften vil gjennem Individet, er det sande
Gode; den, der tænker det Gode, kan ikke Andet end ville
det; at gjøre sig selv til det Almene er intellektuel Kjærlighed
til det Gode, kun i en saadan Kjærlighed er der Selvforstaaelse.
Richards Brøde er nu den, at han ikke theoretisk
vil gjøre sit Selv til det Almene; han elsker
Intet, Ingen, trods sin uhyre Egenkjærlighed ikke engang
sig selv: derfra kommer det Onde. Hans ubøielige
Trods mod Sædelighedens guddommelige Orden er
forenet med en Smerte, en hemmelig Angst, en urolig
Higen, som han i sin Sjæls Bitterhed ikke forstaaer,
fordi han ikke vil forstaae den. De tre ulykkelige Kvinder,
Dronning Margrethe, Hertuginden hans Moder og
Dronning Elisabeth, glemme for et Øieblik det gjensidige
Nag og forene sig som tre Erynnier, for at ud%
% --- p. 540 ---
\setcounter{footnote}{0}%
\opage{540}gyde en Strøm af Forbandelser over Richard, som just
gaaer dem forbi; men han vil Intet høre. »Giv Lyd
Trompeter! --- Trommer slaa Alarm! Lad Himlen ei
de Sladderkvinder høre, som skælde Herrens Salvede;
slaa, slaa!« Som han ved Trods og Trommeslag døver
Forbandelsens Raab, der lyde udenfra, saaledes søger
han ved kløgtigt, sofistisk Ræsonnement at værge sig
mod Forbandelserne indenfra og døve de Erynnier, han
huser i sit Indre. Med al sin gjennemtrængende Forstands
Klarhed er Richard ikke i Stand til at sprede sit
sædelige Mørke; hans Selvforstaaelse er en fortsat Selvbedøvelse:
hans dybere fornuftige Selv kommer ikke for
Lyset.

Slige oversvævende Almenbetragtninger have for den
sædelige Følelse noget Indsmigrende; kun Skade, at de
staae i Strid med Principet. For at hævde det Godes
Ide, vil den atheistiske Ethiker stride med Richard, som
Herkules stred med Kæmpen Antæus, vil løfte ham op
fra Existensens moderlige Jord og kvæle ham i den almene
Tænknings luftige Region; men han forløfter sig
paa Karakteren; Richard er Ethikeren for tung: det er
netop Principet, der tynger. I det abstrakt almene Element
kan der tænkes, men ikke villes, endnu mindre
handles. Almenmennesket kan ikke handle; thi Almenmennesket
er et Begreb, og et Begreb kan ikke handle.
For at bestemme det Virkelige, maa det Almene gaae
igjennem det Individuelle; men naar det Almene gaaer
igjennem det Individuelle, saa gaaer Nødvendigheden ind
under Friheden, Tanken ind under Villien, og saa har
Richard et sikkert Fodfæste. Allerede den Bebreidelse,
at der er en indre Stemme, han overdøver, en Tanke,
han ikke tænker, fordi han ikke vil tænke den, tager
% --- p. 541 ---
\setcounter{footnote}{0}%
\opage{541}han til Indtægt; thi herved er jo uden videre indrømmet,
at de sædelige Karaktertanker ere Villiestanker,
at Tankernes Gang beroer paa Villien. Med denne Indrømmelse
er Richard tilfreds; hans Tanke er i alle hans
Handlinger sig selv lig og gaaer konsekvent i Eet med
hans Villie; Selvvillien er hans Lov, den eneste Lov,
han ikke krænker: at ville i Alt sin egen Villie er hans
Ære. Men paa disse Vilkaar er han endog i bogstavelig
Overensstemmelse med Ethikeren.

»Der gives Noget,« siger vor Ethiker, »som gjælder
mig over Alt, og for hvilket jeg sætter alt Andet
tilside, fra hvis Hævdelse jeg ikke ved nogen mulig
Følge lader mig afholde, Noget, for hvilket jeg uden
Betænkning vilde opoffre min hele jordiske Velfærd, mit
gode Navn og Rygte, mit Liv, den ganske Verdens Vel,
om det kunde komme i Strid dermed. Jeg vil kalde
det Ære.« Denne Fichtiske Haandfæstning vil Richard
strax uden Vægring underskrive; ja, han vil ordenlig
blive rørt og af Henrykkelse falde Ethikeren i Armene,
naar han hører den nærmere Forklaring. »Denne Ære
sætter jeg ingenlunde i Andres Dom over mine Handlinger,
end ikke, om det kunde være min Tidsalders og
Efterverdenens enstemmige Dom; men i den, jeg selv
fælder over den.« Det klinger! Richard lytter opmærksomt
til og hører videre: »Den Dom, jeg fælder over mine
Handlinger, afhænger af, om jeg ved dem bliver i Overenstemmelse   % PRINTED AS IS (p. 541): „Over-|enstemmelse“ with one s
med mig selv, eller ved dem kommer i
Modsigelse med mig selv. I første Tilfælde kan jeg
billige dem; i sidste Tilfælde vilde jeg blive vanæret
for mig selv«\dots{} »Konsekvens, Overensstemmelse med
sig selv!« hvor det klinger i Richards Øren! til at
% --- p. 542 ---
\setcounter{footnote}{0}%
\opage{542}overdøve saa liflige Stemmer ønsker han vist ingen
Trommeslager. Det Almenfornuftige, hvoraf den atheistiske
Ethiker gjør saameget Væsen, kan da heller ikke
volde ham Bryderier. Mine Beslutninger ere ethiske,
naar de ere saaledes beskafne, »at jeg med mit bedste
Vidende kan alvorlig ville, at alle fornuftige Væsener i
samme Stilling skulle fatte de samme Beslutninger«.
Det er jo et lovformeligt Forsvar for Richard. Enhver,
der tænker sig at være af samme Natur som han, i en
med Forbrydelser opfyldt Tid sat med hans Kræfter,
paa hans Plads, i hans Sted, maa jo sige sig selv, at
han i Grunden maatte handle som han. Naar der i et
borgerligt Samfund er Krig af Alle mod Alle, hvad gjælder
da uden den Stærkeres Ret? I denne Henseende
gjør Richard kun, hvad Fornuften byder: han slaaer sine
Fjender til Jorden, fordi han erkjender, at Enhver af
hans Fjender, naar han var i hans Sted, fornuftigvis
maatte slaae ham til Jorden: han handler konsekvent.
Konsekvensen af hans Handlemaade medfører ingenlunde
Samfundets Ødelæggelse; den fører kun til, at den Stærkere
beseirer de Svage, indtil de stridende Kræfter
komme i Ligevægt. Hykleri, Bagtalelse, Mened og Mord
ere Midler, der kun bør anvendes, hvor de behøves; de
falde bort af sig selv, naar Tiden er bleven en anden.
Hvor konsekvent Richard er, kan sees af hans Forhold
til Buckingham. Denne Tyranniets Haandlanger, der
hidtil har fulgt Richard som Ven og Raadgiver, faaer
pludselig Skrupler og vil ikke lade sig bruge til at myrde
de unge Prindser; det er en Inkonsekvens: ved at komme
i Modsigelse med Richard kommer Buckingham i Modsigelse
med sig selv: intet Under, at en saadan Mod%
% --- p. 543 ---
\setcounter{footnote}{0}%
\opage{543}sigelse koster ham hans Hoved. Richard derimod er,
hvad Principet fordrer, i Et og Alt konsekvent, saa
ubøielig konsekvent, at ingen Stoiker med sin »Selvoverenstemmelse«
til Feltraab, skal gjøre ham Rangen
stridig. Hans Hjertes Haardhed beviser just, at han,
ophøiet over sandselig Medlidenhed, kun føler sig »som
aandig Magt«. Han nærmer sig virkelig den Theori,
ifølge hvilken Selvets Tænkning skal med Videns absolut
objektive Nødvendighed kunne slaae over i Villie, i
Beslutning, i Handling. Han handler ikke af blot Naturtilskyndelse
eller af bar Ondskab, men efter en strengt
gjennemtænkt og dybt anlagt Plan; idet han handler med
Frihed, gjør han kun, hvad der er nødvendigt. Henrik
den Sjette myrdes, det er nødvendigt; Clarence myrdes,
det er nødvendigt; Edvards Sønner myrdes, det er aldeles
nødvendigt; Buckinghams Hoved falder, ogsaa det
er nødvendigt. Og hvad er saa Friheden? Just, hvad
en atheistisk Metafysik saa sindrigt forklarer os, Nødvendighedens
udviklede Begreb, dens immanente Overgang
i fornuftig Selvbevidsthed! Og saa skulde en Helt
som Richard, en Karakter, der i saa høi Grad tilfredsstiller
Theorien, ikke kunne siges at forstaae sig selv,
om end hans Selvforstaaelse er i det Onde?

For den atheistiske Ethiker er der da kun den Udvei
tilbage, at opgive det blot Formelle, lægge Eftertrykket
paa Sædelighedens Indhold og appellere til Ideen.
Hvad der stemmer med Ideen er godt, hvad der strider
imod den er ondt; det er ikke den konsekvente Handlemaade
i og for sig, men Karakterens Stilling til Ideen
som Ide, der i Ethiken gjør Udslaget. Richards Forblindelse
viser sig saa deri, at han krampagtig fastholder
% --- p. 544 ---
\setcounter{footnote}{0}%
\opage{544}sit egoistiske Selv og krænker Ideen under enhver Skikkelse;
Venskab, Familie, Stat, Alt offrer han for sin
Egoisme. Han er værre end en Macchiavelist; thi
Macchiavelisten sætter dog Ideen i Statens Form som
det Absolute, men Kong Richard ændser ikke det almene
Vel, han tager Staten som et Rov: for ham er Intet
absolut uden hans eget herskesyge Jeg. Dette Brud
med alle sædelige Magter, er den dybe Modsigelse i
Richards Væsen, og Konsekvensen af dette Brud et
Gjennembrud, der ender med Fortvivlelse. Efter Thronbestigelsen,
endnu før Rygtet melder ham, at Henrik
af Richmond er landet med en Hær, begynder det at
gaae tilbage med Richards Selvtillid. Hans Kraft er
brudt, hans Sjæl har ingen Fred, han er søvnløs; Søvnen
selv, der ellers vederkvæger den Trætte, har, naar
den endelig kommer, ingen Vederkvægelse for ham, den
martrer ham med tunge Drømme; nagende Erindringer
i Strømme af Blod, Stemmer fra Gravene, der raabe: Hævn!
Broderbaandets Sønderrivelse, Moderens Forbandelser,
fuld Forvisning om Alles Fjendskab, Alles Afsky, Alles
Had, dette tilsammen knuger Richards Sind, kuer hans
Mod og lader ham føle, at der er en sædelig Magt, der
straffer den Onde. Saa nærmer sig den sidste Katastrophe.
Richard sover i sit Telt paa »Bosworth
Mark« og drømmer. De Myrdedes Aander vise sig for
ham, En efter En, hver tynger paa ham med sit
Minde, sin Klage og gjennem alle Minder, alle Klager
lyder Dommens Ord: »fortvivl og dø!« Forfærdet farer
Richard op, --- »det er døde Midnat; de kolde Draaber
boble paa hans Legeme«. Saa følger den berømte Monolog,
hvori han for første og da ogsaa for sidste Gang
% --- p. 545 ---
\setcounter{footnote}{0}%
\opage{545}faaer sit ethiske Selv i Tale. Nu endelig forstaaer han
sig selv, men Selvforstaaelsen er den forfærdeligste
Selvmodsigelse. »Hvad frygter jeg? Mig selv? jeg er
alene: ja Richard elsker Richard --- jeg er jeg. Er her
en Morder? Nei; --- jo jeg er en, da fly! --- Hvad, fra
mig selv? God Grund: hvorfor? jeg hævner. Hvem?
mig, paa mig selv? jeg elsker jo mit eget Selv; hvorfor?
for alt det Gode, jeg har gjort mig selv. O, nei!
Ak nei! jeg hade bør mig selv for hadefulde Idræt mod
mit Selv.« Den ethiske Selvmodsigelse er en bogstavelig
Selvfordømmelse; for at faae Richard dømt, behøve vi ikke
at appellere til en oververdslig Ret; Retten er allerede sat:
det er ingen alvidende Gud paa Dommens Dag, men
slet og ret den sædelige Ide, der dømmer Richard.
Omgiven af Dommens Rædsler staaer han som forvandlet.
Da han begyndte med at erklære, at han vilde
være Nidding, da var det kun en overmodig Spot; der
var i hans Bevidsthed ingen sædelig Forskjel mellem en
Nidding og en ædel Helt; men nu, da han er ramt af
Aandens Lyn og i sin Sjæleangst udbryder: »Jeg er en
Skurk!« nu er Forskjellen mellem det Onde og det
Gode gaaet op for ham: Ideen er kommen til Aabenbarelse,
men for at knuse ham. Hvor sønderknust han
er, høre vi af hans Fortvivlelse, hans sjælerystende
Selvbekjendelse; thi »hans Samvittighed har tusind Tunger,
hver Tunge bringer frem en særskilt Klage, hver
Klage ham fordømmer, som en Nidding. Mened og falskest
Ed i største Grad, Mord, grusomt Mord i frygteligste
Grad, og alle Synder i hver enkelt Grad, til Skranken
trænge sig og raabe: Skyldig, Skyldig!«
Dersom det nu virkelig forholdt sig saaledes, at det
% --- p. 546 ---
\setcounter{footnote}{0}%
\opage{546}Godes Ide alene ved sig selv uden Grund i Guds Villie
har faaet Richard dømt og gjort Retfærdigheden Fyldest,
da har den atheistiske Ethik ved denne immanente
Domsakt unægtelig feiret en stor Triumf. Det er da
een Gang for alle gjort vitterligt, at Forskjellen mellem
Godt og Ondt ikke alene kan bestaae, men endog først
da blive objektivt begrundet, naar Gud er fornægtet, og
den uklare Forestilling om en hinsidig Dommer fuldstændig
opgivet. Men er der i nogen videnskabelig
Undersøgelse Grund til at tale om »Tilsnigelse«, »Eklekticisme«,
»Selvmodsigelse« o. desl., da er det nok her.
Den hele Begrundelse beroer paa en Tilsnigelse, idet
man umærkelig lader den \emph{sædelige}, umiddelbart sædelige
\emph{Følelse} træde i Stedet for det \emph{sædelige Princip.}
At den sædelige Følelse saares, naar det Onde sættes ligt
med det Gode, er vist, ligesom det ogsaa er vist, at
Digteren har poetisk Ret, idet han lader Richards bedre
Natur komme til Gjennembrud og bryde hans Trods.
Men det er jo slet ikke det, hvorom her handles. Spørgsmaalet,
det Spørgsmaal, Atheisten her skulde besvare,
angaaer ikke Følelsen, men Principet, det Princip, i
Kraft af hvilket Modsætningen mellem Godt og Ondt
alene har Gyldighed. Er Richard virkelig dømt i Kraft
af det atheistiske Princip? Ingenlunde! Han er i egenlig
Forstand slet ikke sædelig dømt; han er kun opskræmmet,
overvældet af Drømmesyner; hvad han siger, siger han, om
ikke halvt i Søvne, saa dog i en Tilstand, hvori han ikke
er sig selv og sit Princip mægtig. Saasnart han ret er
vaagnet og kommen til Besindelse, fragaaer han da ogsaa
Alt, hvad han i den første Forskrækkelse har sagt. »Samvittighed,
Du feige, Du forfærder!« Denne Indlednings%
% --- p. 547 ---
\setcounter{footnote}{0}%
\opage{547}replik angiver tydelig nok, hvorhen den følgende Strøm af
Selvbebreidelser vil føre. Richard har havt et svagt
Øieblik, og svage Øieblikke kan jo den Stærkeste have,
enten han forøvrigt er god eller ond. Til Ratcliff som
just er traadt ind i Teltet, siger han, at han har havt
en rædsom Drøm; han »frygter, frygter«! Paa Ratcliffs
opmuntrende Ord: »Min gode Fyrste, vær ei angst for
Skygger«, svarer han: »Nu ved Apostlen Paulus, Nattens
Skygger har slaget Richards Sjæl med mere Rædsel,
end tyve tusind virkelige Stridsmænd, bevæbnede,
skudfrie, under Richmond kunde«. Hvad det er for en
Rædsel, hvormed han er slagen, og hvad det er for en
Frygt, han mener, er klart nok: det er ingenlunde Frygt
for Samvittighedens Dom, men alene Frygt for at tabe
Slaget. Skulde Richard i Principet være sædelig overvunden,
da maatte han jo med fuld Besindelse være bragt til
at give Ideen Ret og at anerkjende Samvittigheden. Men
langtfra. Neppe har han reist sig, faaet sin Rustning
paa og skal begynde Slaget, før vi i hans Kommandostemme
atter høre den vilde Trods: »Gaaer mine Mænd,
enhver til sin Forretning, lad ingen Sladderdrømme
skrække Sjælen; Samvittighed er Lyd, som Feige bruge,
indført til Styrkens Tugt. Vor stærke Arm Samvittigheden
være, Sværd vor Lov; saa styrter frem!« Saa
skal »Ideen« da ikke rose sig af, at den i Samvittighedens
Navn har overvundet Richard!

»Men«, siger man, »just fordi Ideen ikke kan beseire
Richard subjektivt, virker den som objektiv Magt og
beseirer ham objektivt.« Nu har man forladt det Ethiske,
og Misforstaaelsen er i fuld Gang; thi en Pil, en Kugle,
en Sten, en Lynstraale kan jo ogsaa beseire et Individ
% --- p. 548 ---
\setcounter{footnote}{0}%
\opage{548}objektivt: hvilken Ære for Ideen at blive sat paa lige
Trin med en simpel Naturmagt! Ideen, siger man, er
ikke Naturmagt, den er Aandsmagt, Aanden selv i Historien,
i Frihedens Verden; som nærværende, Individer,
Handlinger og Begivenheder iboende Aand behersker
den Alt og Alle. Det er kun i egen Indbildning, at
Richard har kunnet løsrive sig fra Ideen; han er i dens
Vold; den holder ham fast med usynlig Haand fra først
til sidst, den bruger ham som Redskab. Blandt de
Mægtige, som i hine Rædslers Periode alle havde forbrudt
sig mod Samfundet, er Richard, den bloddryppende
Heros, der hugger Hovederne af Borgerkrigens Hydra,
netop ved Ideens Magt den mægtigste. »Han er«, siger
Røtscher, »som Verdensaandens Nemesis. Ved denne Betragtning
vinder Richard allerede et ophøiet Postament.
Hans Grundlag er den oprevne Stat, hans Maal er Kronen;
hans absolute Betydning er at være Bæreren af
den historiske Aand, ved hvem alle de Hoveder afmeies,
der have undergravet Almenvæsenet, opløst Staten og
betynget sig med Forbrydelser. Idet han nivellerer Alles
Ret, jævner han Pladsen, at Statens Enhed igjen kan
hæve sig frem af Interessernes Splidagtighed\dots{} Men
som Produkt af den hele, sig opløsende Stat gaaer Richard
selv tilgrunde, thi Forbrydelsen maa sones og
Værktøiet kastes bort, saasnart det har fuldbragt sit
blodige Arbeide«. Fra et historisk-æsthetisk Standpunkt
er der Intet at indvende mod denne Dom; men
hvad bliver det til med Ethiken? Skjøndt Richard
just ikke staaer paa nogen objektiv fortrolig Fod med
Ideen som Ide --- en saadan objektiv, intim Fortrolighed
maa være Stuemetafysikeren forbeholdt --- saa er han
% --- p. 549 ---
\setcounter{footnote}{0}%
\opage{549}jo dog for sin Tid Ideens ivrigste Tjener, dens troeste
Haandlanger, den ufortrødneste Arbeider i dens Vingaard;
objektivt forstaaet, er Aanden da ogsaa langt
mere tilfreds med ham end selv med den fromme Kong
Henrik, for ikke at tale om de Andre. Havde Richard
været en Helgen eller dog en ædel, retskaffen Mand,
saa kunde Historiens Aand jo slet ikke have brugt ham
til sit blodige Værk. Men naar Richard kun har gjort,
hvad Ideen fordrede: efter hvilken Lov kan Ideen da
fordømme ham? Det hændes vel undertiden, at en Magthaver
bortkaster det Redskab, der mente at fortjene
Tak for at have opfyldt sin Herres Ønske. Saaledes
siger Bolingbroke (Richard II. 5 Akt, Slutn.) til Exton,
der bringer ham Richard den Andens Lig i Kisten:
»Vent ingen Tak af mig! thi Du har lagt med blodig
Haand en Skændsel uden Lige paa mig og mit berømmelige
Rige«. Exton minder Kongen om hans utvetydige
Opfordring: »Dit Ord, min Herre, voldte min Bedrift«,
hvortil Bolingbroke (Henrik IV) svarer: »Den elsker
ikke Gift, som bruger Gift; jeg ikke Dig; hans Død var
mig fornøden; men Dig jeg hader, elsker ham i Døden;
Samvittighedens Nag skal lønne Dig, men ingen Ros og
Fyrstegunst af mig; gaa fredløs om som Kain i Nattens
Gys, og vis Dit Hoved ei for Dagens Lys!« Ædel Følelse,
høi Pathos, et herligt Sprog; Poesien jubler, men
--- Ethiken græder.

Skulde Ideen virkelig kunne behandle Richard, som
Bolingbroke behandlede Exton, da vilde det see sort ud
med Adskillelsen af Godt og Ondt. Hvem bærer vel
Skylden, »Giften« eller »den, der bruger Giften?« --- »Forbrydelsen
maa sones!« dog vel ikke den Forbrydelse, at
% --- p. 550 ---
\setcounter{footnote}{0}%
\opage{550}have tilfredsstillet Ideen? Naar der fra atheistisk Side
tales om Richards Forbrydelser, saa tales der slet og
ret subjektivt ud af naturlig menneskelig Følelse, men
ingenlunde metafysisk objektivt ud af »Historiens Aand«.
Det er i det Hele betegnende for de gudsfornægtende
Theorier under alle Former, at man, efter at have sat
det Almenfornuftige i Guds Sted, nødes til at personificere
Ideen og, for at bevare et Sædelighedens Skin,
tillægge den en Art Samvittighed og behandle den som
moralsk Personlighed. Man pukker paa Begrebet og indsmugler
Forestillinger, der staae i skærende Strid med Begrebet;
man forguder Abstraktionen og sminker det døde
Abstraktum op til at see ud som et sædeligt Væsen, der
belønner de Gode og straffer de Onde: det er en for Liv og
Videnskab lige uhyggelig Mystifikation. Ideen, den fra
Guds hellige, retfærdige Villie løsrevne Ide, er ingen
sædelig Magt, hverken god eller ond, hverken retfærdig
eller uretfærdig. Dens Væsen er i Naturens Orden, i
Kræfternes Virken, i Grundenes Nødvendighed: den lever
i objektive Loves Opfyldelse. Dens Ret er en aandig
Naturret: Alt, hvad der er muligt, er tilladt, kun det
Umulige er forbudt. Ogsaa det Onde skeer efter Love;
derfor er Ideen ligesaa forsonet med det Onde som med
det Gode; den er uden Glæde, men ogsaa uden Smerte,
uden Vrede, men ogsaa uden Medlidenhed; den harmes
ikke paa Vold og Rov, den sørger ikke, naar en Bøddel
krones, saalidt som Øxen sukker, fordi en Uskyldig
henrettes.

Resultatet af vor Undersøgelse er dette: paa Atheismens
Standpunkt gives der en ligesaa klar og gjennemgaaende
Selvforstaaelse i det Onde som i det Gode, og
% --- p. 551 ---
\setcounter{footnote}{0}%
\opage{551}det af den simple Grund, at det atheistiske Princip ikke
formaaer at opretholde den sædelige Forskjel mellem
Godt og Ondt. Det nytter ikke at beraabe sig paa
Menneskets Natur, thi Menneskenaturen er kun til i virkelige
Mennesker, og de virkelige Mennesker have hver sin
Natur, forsaavidt de have hver sit Sind. Den, der har
et godt Sind, foretrækker det Gode, den, der har et ondt
Sind, det Onde; skal Sædeligheden umiddelbart beroe
paa Sindet, da er den Onde ligesaa berettiget som den
Gode. Det nytter heller ikke at beraabe sig paa Menneskets
fornuftige Væsen, thi Tanken, der skal opfatte det
Almenfornuftige, maa, --- naar den da ikke ved tomme
Theorier vil udelukke sig selv fra Liv og Virkelighed, ---
nødvendigvis sætte det fornuftige Væsen i Menneskets
tænkende og villende Selv og gjøre den konsekvente
Selvtilfredsstillelse til øverste Grundsætning. Men naar
Egenvillien er Lov, og Konsekvensen Maalestok, kan den
Onde paa sin Vis have en ligesaa konsekvent Selvtilfredsstillelse
af at gjennemføre sin Villie, som den Gode af at
gjennemføre sin. Tilfredsstillelsen er paa ingen af Siderne
fuldkommen; den Onde har en Modstand i sin bedre Natur
og derved i sig selv; ligesaa har den Gode en Modstand i sin
lavere Natur og derved i sig selv; Vilkaaret for denne som
for hin er uundgaaelig Kamp og Strid. Hvor stor Tilfredsstillelsen
vil blive, beroer alene paa den Karakterstyrke,
hvormed Principet gjennemføres. Den konsekvente Karakter
har en stor, den inkonsekvente en ringe Selvtilfredsstillelse, uden alt Hensyn paa Godt og Ondt. Efter
at have gjort Principet saa subjektivt som muligt og sat
det Almene i Selvvillien, er det forgjæves at ville føre
Sagen ind for et objektivt Forum og lade Ideen dømme.
% --- p. 552 ---
\setcounter{footnote}{0}%
\opage{552}Den sædelige Ide er Villieside; har den ikke sit Midtpunkt
i Guds Villie, da er den ubetinget afhængig af
Menneskets Villie; Ideens Dom over Mennesket er da
kun Menneskets egen Dom over sig selv: Dommen er
ubetinget subjektiv. Skal Ideen ikke desmindre hæve
sig over det Subjektive og gjælde for selvgyldig objektiv
Magt, da ophører den at være ethisk.

Herved er ingenlunde bevist, at en Atheist ikke
skulde kunne være en sædelig Karakter, tvertimod! det
er netop bevist, at han, eftersom hans Sind og Subjektivitet
er, kan være baade sædelig og usædelig, idet det
er vist, at Modsætningen af Sædeligt og Usædeligt, ikke
for den enkelte Atheist, men for det atheistiske Princip
er uden Betydning.

Der er, som bekjendt, i Konst, Poesi og Historie
visse store orienterende Synspunkter, om hvilke den
Dannede, hvorlidet han iøvrigt har Leilighed til at fordybe
sig i noget af de nævnte Fag, ikke let vil tilgive
sig selv at være ganske uvidende. Burde det
Samme nu ikke ogsaa gjælde, naar Talen er om Religion
og Sædelighed? Blandt de i denne Henseende orienterende
Hovedsætninger henregne vi følgende:

Guds Tilværelse, Samvittighedens Ret og Sjælens
Udødelighed ere personlige Forudsætninger, hørende saa
ubetinget til Sjælens indre Liv, at deres Gyldighed, naar
vi tage det strengt, ligesaa lidt kan bevises som modbevises
af Videnskaben. De ere Troessætninger, hvis
Eiendommelighed Videnskaben vel kan belyse, men ikke
begrunde.

Theologien, der vil begrunde, og Atheismen, der
vil undergrave den aabenbarede Religion, have begge lige
% --- p. 553 ---
\setcounter{footnote}{0}%
\opage{553}megen Uret. Atheismens umiskjendelige Overlegenhed
beroer alene derpaa, at Striden føres i Videns Element,
et Element, der kvæler Troen. Theologiens Afmagt er
intet Bevis mod Religionen, thi Religionen med sit Princip
har Intet med Theologien at skaffe.

Det sande personlige Bevis for Sædelighedens Gyldighed
er uadskilligt fra det sande personlige Bevis for
Religionens Gyldighed; Beviset støtter sig til Villiens
ubetingede Realitet. Den menneskelige Villie mister
øieblikkelig sit sande Evighedspræg, sin absolute Afhængighed
og dermed tillige sin høiere sædelige Selvstændighed,
naar Troen paa Guds Villie er udelukket.

Først naar det religiøs-sædelige Princip er anerkjendt
i fuldkommen Uafhængighed af objektiv Viden og
Videnskab, kan det med alsidig Konsekvens udfolde sine
Erkjendelsers Indhold i sin egen Sfære og Religionslæren
opstaae gjenfødt til fornyet Virksomhed. At ville fra
et atheistisk Standpunkt angribe denne Synsmaade ved
at slaae om sig med »videnskabelige Modsigelser«, er
Nonsens: da Videnselementet ifølge vort valgte Standpunkt
er Religionens Væsen imod, saa følger det jo
ganske af sig selv, at Religionslæren, naar den behandles
i Videnskaben, nødvendigvis maa vrimle af Modsigelser.
Den eneste Vei til at blive Modsigelserne kvit, er
at skille Religionen fra Videnskaben.{\renewcommand{\thefootnote}{*)}\footnote[1]{Jvfr. Religionsfilosofi S. 5—11, 74---81, 97—108 ff.}}

Ved at fastholde disse Grundsætninger ville dannede
Læsere let kunne orientere sig i det Virvar af theologiske
og atheistiske Forsikkringer, der sættes i Omløb
under Paaberaabelse af »Grunde«, »klare Grunde«,
% --- p. 554 ---
\setcounter{footnote}{0}%
\opage{554}»uigjendrivelige Grunde«. Et er det at raabe med
Grunde, et Andet at give Grunde; det Første er lettere,
men det Sidste er paalideligere. De, der ivrigst raabe
med Grunde, ere ikke altid paalidelige, ikke saa paalidelige,
som de, der raabe med Reier: de Sidste have
dog i Reglen den Vare, de falbyde, men de Første have
desværre tit ikke Andet end Raabet.

\signature{R. \emph{Nielsen}.}

\end{document}
